\documentclass[varwidth=14cm,border=10pt]{standalone}
\usepackage{amsmath,xcolor,array}
\usepackage{fontspec,unicode-math}
\setmainfont{DejaVu Sans}[ItalicFont={DejaVu Sans},ItalicFeatures={FakeSlant=0.2},BoldItalicFont={DejaVu Sans Bold},BoldItalicFeatures={FakeSlant=0.2}]
\setmathfont{texgyredejavu-math.otf}
\setmathfont{DejaVu Sans}[range={up,bfup}]
\setmathfont{DejaVu Sans}[range={it,bfit},FakeSlant=0.2]
\setmathfont{texgyredejavu-math.otf}[range={"2202,"1D715,"03B5,"03F5,"2112,"2131}]
\pagecolor[HTML]{1E1E2E}\color[HTML]{F4F4F4}
\definecolor{accent}{HTML}{89B4FA}
\definecolor{warn}{HTML}{F38BA8}
\definecolor{good}{HTML}{A6E3A1}
\newcommand{\heading}[1]{\par\medskip{\large\bfseries\color{accent}#1}\par\smallskip}
\setlength{\parindent}{0pt}
\raggedright
\begin{document}

{\Large\bfseries A2 Q1 — where $K$ and $P$ come from}

\heading{1. The one algebraic identity behind it}
A spring of stiffness $k$ joining carts $i$ and $j$ stores
\[
  \tfrac12 k\,(x_j-x_i)^2
  \;=\;
  \tfrac12 k\,x_i^2+\tfrac12 k\,x_j^2-k\,x_ix_j .
\]
Read off the quadratic form $\tfrac12x^{\!\top}Kx$: that single spring puts
\[
  +k \text{ on } K_{ii}, \qquad +k \text{ on } K_{jj}, \qquad -k \text{ on } K_{ij}=K_{ji}.
\]
A {\bfseries grounded} spring is the same with the wall at $x=0$:
$\tfrac12k\,x_i^2$, so it adds $+k$ to $K_{ii}$ only --- no off-diagonal,
because a wall is not a design variable.

\heading{2. The rule that falls out}
\[
  \boxed{\;K_{ii}=\!\!\sum_{\text{springs touching } i}\!\! k,
  \qquad
  K_{ij}=-\,k_{\text{(the spring joining } i,j)}\;}
\]

\heading{3. Your eight springs, added up}
\[
  \begin{array}{@{}l@{}}
    K_{11}=k_1+k_4+k_5=5000+1000+2500=8500\\[1pt]
    K_{22}=k_2+k_4+k_6=1500+1000+500=3000\\[1pt]
    K_{33}=k_3+k_5+k_6+k_7+k_8=2000+2500+500+3000+3500=11500\\[4pt]
    K_{12}=-k_4=-1000,\qquad K_{13}=-k_5=-2500,\qquad K_{23}=-k_6=-500
  \end{array}
\]
Cart 3 has the big diagonal because {\bfseries five} springs touch it: ground,
two couplings, and two to the right wall.

\heading{4. Where $P$ comes from}
The loads do {\itshape work}, which {\itshape lowers} energy:
\[
  -1000x_1-2000x_2-3000x_3 \;=\; -\,x^{\!\top}P,
  \qquad P=(1000,2000,3000)^{\!\top}\text{ N}.
\]
That is the only reason the minus sign is there, and the only reason the problem
has a finite minimum --- without $P$, $V=\tfrac12x^{\!\top}Kx$ is minimised by
$x=0$.

\heading{5. Why this form is worth collecting}
\[
  \nabla V = Kx-P \;=\; 0 \quad\Longleftrightarrow\quad \boxed{\,Kx=P\,},
  \qquad \nabla^2V = K .
\]
{\color{good}Minimising the potential energy \emph{is} solving equilibrium.}
The optimisation problem and the FE problem are the same problem --- which is
exactly what part 4 checks by assembling $K$ a second way, from
$k_e\!\left[\begin{smallmatrix}1&-1\\-1&1\end{smallmatrix}\right]$ elements, and
getting the same matrix.

\heading{6. Free checks on your $K$}
\begin{itemize}\itemsep0pt
\item Symmetric: $K=K^{\!\top}$. Springs have no preferred direction.
\item Diagonal dominance: each row's diagonal $\ge$ sum of $|$off-diagonals$|$,
      strictly so wherever a grounded spring sits.
\item Positive definite: eigenvalues $2714.4,\ 7369.8,\ 12915.8>0$, so the
      stationary point is a {\bfseries minimum}, and the only one.
\item {\color{warn}Rows do not sum to zero here} --- they would for a free
      floating structure, but grounding to the walls is what removes the rigid
      translation and makes $K$ invertible.
\end{itemize}

\heading{7. Yes --- this is Lecture 3, slides 30--38}
L03 does the identical thing on a smaller system (Rao Fig.~6.21): two bodies,
four springs, loads 1000 and 2000~lb.
\[
  V(x_1,x_2)=\tfrac12(300)x_1^2+\tfrac12(400)(x_2-x_1)^2+\tfrac12(200)x_2^2
             +\tfrac12(100)x_2^2-1000x_1-2000x_2
\]
and the general statement it gives you is
\[
  V(x)=\underbrace{\sum_i \tfrac12 k_i\Delta_i^{\,2}}_{\text{strain energy}}
       -\underbrace{\sum_j P_j x_j}_{\text{work of the loads}} .
\]
The slide titled {\itshape The two routes meet} is this note in three lines:
\begin{enumerate}\itemsep0pt
\item Expand and collect: $V=\tfrac12x^{\!\top}Kx-x^{\!\top}P$, {\itshape with
      exactly the assembled $K$}.
\item $\nabla V=Kx-P=(\text{internal force})-(\text{applied load})$ --- the
      out-of-balance force.
\item $H=K$.
\end{enumerate}
And the punchline A2 part 4 is checking:
\[
  x^{(2)}=x^{(1)}-H^{-1}\nabla V = x^{(1)}-K^{-1}\!\left(Kx^{(1)}-P\right)=K^{-1}P .
\]
{\color{good}The starting point cancels: one Newton step \emph{is} the assembly
solve.} {\color{warn}Only because the springs are linear} --- add large
deflections, plasticity or contact and $H$ is no longer constant, which is when
a real FE code must iterate.

\end{document}
